% Generated draft: requires mathematical and source review.
\documentclass[11pt]{article}
\usepackage[margin=25mm]{geometry}
\usepackage{fontspec}
\setmainfont{Latin Modern Roman}
\usepackage{amsmath,amssymb,mathtools,unicode-math}
\setmathfont{Latin Modern Math}
\usepackage{enumitem,needspace,xurl,hyperref}
\setlength{\parindent}{0pt}
\setlength{\parskip}{5pt}
\setlength{\emergencystretch}{3em}
\newcommand{\N}{\mathbb{N}}
\newcommand{\Z}{\mathbb{Z}}
\newcommand{\Q}{\mathbb{Q}}
\newcommand{\R}{\mathbb{R}}
\newcommand{\C}{\mathbb{C}}
\newcommand{\F}{\mathbb{F}}
\newcommand{\A}{\mathbb{A}}
\newcommand{\PP}{\mathbb{P}}
\newcommand{\E}{\mathbb{E}}
\newcommand{\eps}{\varepsilon}
\newcommand{\dd}{\,\mathrm{d}}
\begin{document}
\section*{UnsolvedMath ID: 700729}
\textbf{Source problem code:} catalogue entry
\subsection{Definitions and mathematical statement}
Let \(F_m\) be a free group of rank \(m\), and define the free Burnside group
\[
B(m,n)=F_m/\langle\!\langle w^n:w\in F_m\rangle\!\rangle,
\]
where the denominator is the normal closure in \(F_m\). Put \(G=B(m,n)\). The source assumes that \(n\) is an odd positive integer with \(n\gg1\), meaning sufficiently large, but supplies neither a numerical threshold nor its possible dependence on \(m\) or on the part under consideration. The general range of the rank parameter \(m\) is not specified; the additional rank restrictions below are explicit source assumptions.

The following are the questions proposed by S. V. Ivanov in Kourovka Notebook Problem 11.36.
\begin{enumerate}
\item[(a)] Is every subgroup \(H=\langle x,y\rangle\), with \(x,y\in G\), isomorphic to the Burnside \(n\)-product of two cyclic groups? The source does not define this product, specify the orders of its cyclic factors, or address degenerate generating pairs. Thus the product terminology is retained without replacing it by a particular quotient construction.
\item[(b)] Assume \(m>1\). For every automorphism \(\varphi\in\operatorname{Aut}(G)\) satisfying
\[
\varphi^n=\operatorname{id}_G
\quad\text{and}\quad
\varphi(b)\varphi^2(b)\cdots\varphi^n(b)=1
\quad\text{for every }b\in G,
\]
must \(\varphi\) be inner? Here \(\varphi^j\) denotes composition \(j\) times, the product is taken in the displayed order, and being inner means that there exists \(g\in G\) such that \(\varphi(b)=g^{-1}bg\) for every \(b\in G\).

The supplied editors' comment reports affirmative results for large prime \(n\), attributed to E. A. Cherepanov (2006); for prime \(n\ge1009\), attributed to V. S. Atabekyan (2011); and for odd \(n\ge1003\) under the additional hypothesis that the order of \(\varphi\) is a prime power, attributed to V. S. Atabekyan (2014). These are supplied source qualifications, not independently verified assertions.
\item[(c)] Assume \(m<\infty\). Is \(G\) Hopfian, that is, is every surjective homomorphism \(f:G\to G\) injective?
\item[(d)] Is every retract of \(G\) free? A subgroup \(R\le G\) is a retract if there exists a homomorphism \(r:G\to R\) such that \(r(x)=x\) for every \(x\in R\). The word ``free'' is not qualified in the supplied record: whether it means an ordinary free group or a free group in a Burnside variety remains unresolved. These interpretations must not be silently identified.
\item[(e)] This part is explicitly archived as solved in the supplied source. Let \(\Z G\) be the integral group ring, whose elements are finite sums \(u=\sum_{g\in G}u_g g\), with \(u_g\in\Z\), and define
\[
\operatorname{supp}(u)=\{g\in G:u_g\ne0\}.
\]
The original question asks whether
\[
\begin{aligned}
\forall a,b\in\Z G,\quad ab=0\ \Longrightarrow\ 
&\exists a_1,c,d,b_1\in\Z G\ \exists C\le G:\\
&a=a_1c,\qquad b=db_1,\qquad cd=0,\\
&C\text{ is cyclic},\qquad
\operatorname{supp}(c)\cup\operatorname{supp}(d)\subseteq C.
\end{aligned}
\]
This is the source's factorization criterion for trivial zero divisors; in particular it applies to nonzero \(a,b\) with \(ab=0\). The supplied archive reports a negative answer, stating that nontrivial zero divisors exist, and attributes this to S. V. Ivanov and R. Mikhailov (2014). This is an attributed archive report, not a new verification of that result.
\end{enumerate}
\subsection{Short English statement}
For a free Burnside group of sufficiently large odd exponent, the source asks about two-generated subgroups, certain finite-order automorphisms, the Hopf property in finite rank, and retracts. It also preserves an archived question about zero divisors in \(\Z G\), reported as answered negatively. The supplied registry labels the whole record partially solved. The exponent threshold, the Burnside \(n\)-product, and the meaning of ``free'' for retracts require clarification before a fully precise formulation can be accepted.
\subsection{Sources}
\begin{itemize}
\item Supplied local UnsolvedMath registry record: permanent ID \(700729\), source problem code \texttt{catalogue entry}, title \emph{Kourovka Notebook Problem 11.36}; upstream revision \texttt{d371ee917ac9647fbb06cd6fc529deec6f3dda82}. The record is labelled \texttt{partially\_solved}.
\item \emph{The Kourovka Notebook (Unsolved Problems in Group Theory)}, 11th issue (1990), Problem 11.36, proposed by S. V. Ivanov; supplied record describes an update in the 21st issue (2026), arXiv:1401.0300v48. The supplied collection description names E. I. Khukhro and V. D. Mazurov as editors of Issues 1--21 (1965--2026).
\item Editors' comment on part (b): E. A. Cherepanov, \emph{Int. J. Algebra Comput.} \textbf{16} (2006), 839--847, cited for large prime \(n\).
\item Editors' comment on part (b): V. S. Atabekyan, \emph{Izv. Math.} \textbf{75}, no. 2 (2011), 223--237, cited for prime \(n\ge1009\).
\item Editors' comment on part (b): V. S. Atabekyan, \emph{Math. Notes} \textbf{95}, no. 5 (2014), 586--589, cited for odd \(n\ge1003\) when the automorphism has prime-power order.
\item Archive of Solved Problems entry for part (e): S. V. Ivanov and R. Mikhailov, \emph{Canad. Math. Bull.} \textbf{57} (2014), 326--334. The supplied archive reports a negative answer and the existence of nontrivial zero divisors.
\item Supplied status review dated exactly \texttt{2026-10-16}: canonical status \texttt{partially\_solved}, evidence field \texttt{published}; its note places part (e) in the solved archive and the remaining parts in the main section. This is supplied dated curation, not an independently verified literature assessment. No referenced publication was independently consulted for this draft.
\end{itemize}
\end{document}
